\documentclass[11pt,a4paper]{article}
\usepackage[utf8]{inputenc}
\usepackage[T1]{fontenc}
\usepackage[portuguese]{babel}
\usepackage{geometry}
\usepackage{longtable}
\usepackage{array}
\usepackage{tabularx}
\geometry{margin=2.5cm}

\title{RunX -- Relatorio Final Fase 2}
\author{ASW 2025/2026 -- Grupo 16}
\date{15 de abril de 2026}

\begin{document}
\maketitle

\section*{1. Enquadramento evolutivo (Checkpoint 1 $\rightarrow$ Checkpoint 2)}

Este documento e evolutivo em relacao ao checkpoint anterior. Mantem a base de especificacao da fase 1 e descreve a sua concretizacao tecnica no backend da fase 2.

\section*{2. Reflexao por secao do documento base}

\subsection*{2.1 Enquadramento}
O problema e os objetivos iniciais mantiveram-se estaveis. A evolucao principal foi transformar a proposta funcional numa API real com persistencia e regras de negocio.

\subsection*{2.2 Utilizadores e cenarios}
Os papeis Runner e Organizer foram implementados com autorizacao no dominio, sobretudo em grupos e eventos. A experiencia mostrou a importancia de explicitar regras de permissao em servico e nao apenas em controlador.

\subsection*{2.3 Storyboards e fluxos}
Os fluxos de interface foram usados como fonte para contratos HTTP e para o desenho dos endpoints. Nesta fase, a prioridade foi consistencia de request/response e nao detalhe visual.

\subsection*{2.4 Requisitos funcionais}
Os modulos principais ficaram operacionais (auth, users, runs, feed, weekly goals, groups). O reforco mais relevante foi adotar validacao Zod para body, params e query.

\subsection*{2.5 Requisitos nao funcionais}
A arquitetura em camadas e o tratamento centralizado de erro ficaram consolidados. A documentacao OpenAPI foi expandida para todos os endpoints implementados.

\subsection*{2.6 Modelo de dados e API}
O esquema Prisma foi estendido para entidades sociais e colaborativas e alinhado com a API REST documentada no Swagger.

\section*{3. Entregas tecnicas concluidas}

\begin{itemize}
  \item Arquitetura backend em camadas (routes, controllers, services, models, schemas, middlewares).
  \item Persistencia com Prisma ORM e PostgreSQL.
  \item Autenticacao JWT com refresh token.
  \item Modulos de Auth, Users, Runs, Themes, Weekly Goals, Feed e Groups.
  \item Validacao de entradas com Zod em body, params e query.
  \item Documentacao Swagger para todos os endpoints implementados na fase 2.
\end{itemize}

\section*{4. Escopo de integracoes externas}

Nesta fase, integracoes externas (meteorologia em tempo real e sugestao inteligente de rotas) nao sao obrigatorias para a avaliacao funcional base.

Estado atual:

\begin{itemize}
  \item Endpoint de sugestao de rota existe em modo \textit{stub}.
  \item Integracao de meteorologia esta preparada e opcional.
\end{itemize}

\section*{5. Evidencias de validacao}

Foram preparados guioes de testes:

\begin{itemize}
  \item \texttt{docs/auth-tests.md}
  \item \texttt{docs/groups-tests.md}
\end{itemize}

Estado tecnico atual:

\begin{itemize}
  \item Build TypeScript validado com sucesso (\texttt{npm run build}).
  \item API operacional em modo desenvolvimento (\texttt{npm run dev}).
\end{itemize}

\section*{6. Anexo de cumprimento (estimativas)}

\subsection*{6.1 Requisitos gerais}

\begin{tabularx}{\textwidth}{|l|c|X|}
\hline
\textbf{Categoria} & \textbf{Cumprimento} & \textbf{Evidencia} \\
\hline
Arquitetura backend em camadas & 100\% & Estrutura modular aplicada aos fluxos principais da API. \\
\hline
Persistencia relacional & 100\% & Prisma + PostgreSQL com migracoes e entidades de dominio. \\
\hline
Validacao rigorosa de input & 95\% & Zod aplicado em body/params/query nas rotas principais. \\
\hline
Documentacao de API & 100\% & OpenAPI/Swagger completo para endpoints implementados. \\
\hline
\end{tabularx}

\subsection*{6.2 Requisitos especificos por fase}

\begin{tabularx}{\textwidth}{|l|c|X|}
\hline
\textbf{Fase} & \textbf{Cumprimento} & \textbf{Justificacao} \\
\hline
Fase 1 (Especificacao) & 100\% & Documento base completo e preservado como referencia evolutiva. \\
\hline
Fase 2 (Backend) & 95\% & API funcional, compilavel, validada e documentada. \\
\hline
Fase 3 (Frontend + IA) & 30\% & Contratos e stubs preparados; integracao final ainda pendente. \\
\hline
Fase 4 (Entrega final) & 10\% & Preparacao documental e tecnica inicial. \\
\hline
\end{tabularx}

\section*{7. Ajustes finais de documentacao}

\begin{itemize}
  \item README consolidado com setup, endpoints e estado da fase.
  \item Relatorio final consolidado (Markdown e TeX).
  \item Guia de testes para grupos e eventos.
\end{itemize}

\section*{8. Proximos passos}

\begin{itemize}
  \item Integrar funcionalidades IA previstas para fases seguintes.
  \item Cobrir endpoints com testes automatizados.
  \item Alinhar frontend com os contratos finais do backend.
\end{itemize}

\end{document}
